\section{Details of the Correctness Proof}
\label{sec:supp-proof}
This section gives the parts of the proof of
Theorem~\ref{main-thm:encoding-correctness} that the article only outlines
(Section~\ref{main-subsec:proof-outline}); the auxiliary lemmas follow in
Section~\ref{sec:supp-aux}.

\subsection{Semantic Setting}
\label{subsec:semantic-setting}

Let
\[
  w=(a_0,t_0)(a_1,t_1)(a_2,t_2)\cdots
\]
be a timed word, where the timestamps are nonnegative, strictly increasing, and
time-divergent.  We write
\[
  \Delta_i=t_{i+1}-t_i>0,
  \qquad
  \elapsed_w(i,j)=t_j-t_i.
\]
The MTL satisfaction relation is written $w,i\satM f$.

The translation uses an \emph{extended word} $\rho$ containing the same timed
word together with clock values and reset decisions.  In the Coq development,
clocks are not indexed by syntax-tree paths.  The clocks of a formula $f$ are
its timed subformulas:
\[
  \mathit{Clock}(f)=\{F\mid F\text{ is a primitive timed subformula of }f\},
\]
a finite type (\texttt{Clock f}), and $\mathsf{clock\_of}(f,F)$ is the formula
$F$ viewed as an element of $\mathit{Clock}(f)$.  Extended words, LTL formulas
over the extended alphabet, the translation, and the automata are all indexed
by $f$, so every clock they mention is, by its type, a timed subformula of $f$.
Thus, a primitive timed formula $F$ denotes one logical clock independently of
where it occurs in the syntax tree.  If two different paths locate the same
syntactic formula, they therefore locate the same clock.  For a clock
$x\in\mathit{Clock}(f)$ we write
\[
  \CVal_\rho(i,x)\in\Real_{\ge0},
  \qquad
  \CReset_\rho(i,x)\in\{\mathsf{true},\mathsf{false}\}.
\]
We write $\ewbase(\rho)$ for the underlying timed word of $\rho$; the predicate
$\mathsf{same\_base}(\rho,w)$ states that $\ewbase(\rho)=w$.

Clock consistency imposes the recurrence
\begin{equation}
\CVal_\rho(i+1,x)=
\begin{cases}
  \Delta_i,
  & \text{if }\CReset_\rho(i,x)=\mathsf{true},\\[1mm]
  \CVal_\rho(i,x)+\Delta_i,
  & \text{otherwise},
\end{cases}
\label{eq:supp-clock-update}
\end{equation}
together with nonnegativity of all clock values.  Thus, a reset at position $i$
makes the clock measure the time elapsed since $t_i$.  (In the
notation of Section~\ref{main-sec:tba}, $\Delta_i=\delta_{i+1}$.)

The LTL satisfaction relation $\rho,i\satL\varphi$ is the standard one, with
the atoms interpreted on $\rho$: an action $a$ holds iff $a_i=a$;
$\rho,i\satL(x\sim d)$ iff $\CVal_\rho(i,x)\sim d$;
$\rho,i\satL\Reset(x)$ iff $\CReset_\rho(i,x)=\mathsf{true}$; and
$\rho,i\satL\Unch(x)$ iff $\CReset_\rho(i,x)=\mathsf{false}$.

Paths are nevertheless retained for structural traversal.  For a formula $f$
located at path $\pi$, $\T_\pi(f)$ denotes its LTL translation, and
$\Node(\mathit{root},\pi)$ returns the subformula at that path.  The auxiliary function
$\mathsf{clock\_at}(\mathit{root},\pi)$ first locates the formula and then applies
$\mathsf{clock\_of}(\mathit{root},\cdot)$.  Consequently,
\[
\begin{aligned}
  &\Node(\mathit{root},\pi_1)=F=\Node(\mathit{root},\pi_2)\\
  &\quad\Longrightarrow\quad
  \mathsf{clock\_at}(\mathit{root},\pi_1)
  =
  \mathsf{clock\_at}(\mathit{root},\pi_2).
\end{aligned}
\]
This is the content of the Coq lemma
\texttt{identical\_\allowbreak{}nodes\_\allowbreak{}share\_\allowbreak{}clock}.

\subsection{Derived Ordinary Timed Operators}
\label{subsec:derived-operators-proof}

The source notation contains ordinary and hatted timed operators, but only the
hatted ones are primitive in the Coq datatype.  The ordinary operators are the
definitions of Eqs.~\eqref{main-eq:derive-ule}--\eqref{main-eq:derive-rge}, whose
agreement with the standard pointwise semantics is
Theorem~\ref{main-thm:derived-semantic-compatibility}.
The upper-bounded identities separate the case in which the current position
already supplies the relevant endpoint.  For lower bounds, well-formedness
requires $d>0$, so a qualifying endpoint is necessarily strictly later than the
current position; strict upper bounds also require $d>0$.  Since these
equations are Coq definitions (\texttt{MUle}, \texttt{MUlt}, \texttt{MUge},
\texttt{MUgt}, and their Release counterparts), ordinary timed operators are
reduced to Boolean structure plus one hatted timed primitive before the timed
proof cases are considered.  This is why the structural proof has eight timed
branches rather than sixteen.

\subsection{Clock Sharing for Identical Subformulas and Repeated Obligations}
\label{subsec:clock-reuse}

The Coq development establishes two levels of sharing.

\runinhead{Sharing between different syntax-tree occurrences.}
As explained in the semantic setting, two occurrences of the same primitive
timed subformula share one clock even at different paths; this is built into
the semantics of the verified translation, not a post-processing optimization.
In addition, the lemma
\texttt{identical\_\allowbreak{}nodes\_\allowbreak{}share\_\allowbreak{}canonical\_\allowbreak{}trace}
shows that the two occurrences have the same canonical value/reset trace.

\runinhead{Sharing between repeated run-time activations.}
A single timed subformula may also be activated repeatedly under
$\Box$, $\Until$, or $\Release$, so several obligations can overlap even
though they use one shared clock.  The canonical reset policy resolves this
competition by keeping an old reference when it dominates newer activations
and restarting the clock only in the dual cases.

For upper-bounded hatted Until, suppose an older activation started at $i_0$
and a newer one at $i_1$, with $i_0\le i_1<j$.  If $j$ is a valid endpoint for
the older obligation, then
\[
  t_j-t_{i_1}\le t_j-t_{i_0}\le d
\]
and the prefix required by the newer activation is a suffix of the older
prefix.  Therefore keeping the oldest still-useful clock reference cannot lose
the newer obligation.  In the mechanized proof this reasoning is carried by
the residual predicate used in \texttt{complete\_\allowbreak{}UhatLe}.

For lower-bounded hatted Release, an older unresolved reference has a clock
value at least as large as any newer one and therefore reaches the lower
threshold no later.  The residual predicate
$\mathsf{RRes}$ (Section~\ref{subsec:residuals}) records precisely the part of the Release obligation that
remains after some clock age has already accumulated.  Its one-step
preservation property is isolated in the Coq lemma
\texttt{rge\_\allowbreak{}residual\_\allowbreak{}shift} and used by
\texttt{complete\_\allowbreak{}RhatGe}.

The dual families use the opposite policy.  For
$\hUntil_{\ge d}$ and $\hRelease_{\le d}$, a newly active obligation may reset
the shared clock.  The stronger LTL continuation (the continuation $\Gamma$ of
lower-bounded Until, Section~\ref{subsec:clocked-ltl-proof}, and the
weak-until branch of upper-bounded Release) ensures that older semantic
obligations remain represented after the restart.
Each of these arguments holds for the strict comparison of the same family,
with the strict residual predicates of Section~\ref{subsec:residuals}.

\subsection{Translation of the Primitive Timed Operators}
\label{subsec:clocked-ltl-proof}

Consider a primitive hatted timed formula $F$ at syntax-tree path $\pi$, with
left and right subformulas $p$ and $q$.  Let
\[
  x=\mathsf{clock\_of}(F),\qquad
  \overline{P}=\T_{\pi0}(p),\qquad
  \overline{Q}=\T_{\pi1}(q),
\]
and write
\[
  H=\Unch(x),\qquad R=\Reset(x),\qquad
  \alpha\WeakUntil\beta\equiv(\alpha\Until\beta)\lor\Always\alpha.
\]
Let $\precsim\in\{\le,<\}$ and $\succsim\in\{\ge,>\}$, with the complementary
comparisons $\precsim^{\mathrm c}$ and $\succsim^{\mathrm c}$ of
Section~\ref{main-sec:mtlbuchi}.  For lower-bounded hatted Until define
\[
  \Theta=(x\succsim d)\land(\overline{P}\Until \overline{Q}),\qquad
  \Gamma=(\overline{P}\Until \Theta)\lor(\Always \overline{P}\land\Always\Eventually \overline{Q}),
\]
and for the Release cases define
\[
\begin{aligned}
  D&=(x\mathrel{\precsim^{\mathrm c}}d)\lor(\overline{P}\land \overline{Q}),\qquad
  K=(x\mathrel{\succsim^{\mathrm c}}d)\land H,\\
  \Lambda&=(\overline{P}\Release \overline{Q})\lor((x\mathrel{\succsim^{\mathrm c}}d)\land \overline{P}).
\end{aligned}
\]

The complete primitive timed translation mechanized in Coq consists of eight
clauses, of four shapes:
\begin{align*}
\T_\pi(p\hUntil_{\precsim d}q)
  &= \Next\bigl(((x\precsim d)\land H\land \overline{P})\\
  &\qquad\quad\Until((x\precsim d)\land \overline{Q})\bigr),\\[1mm]
\T_\pi(p\hUntil_{\succsim d}q)
  &= R\land\Next\Gamma,\\[1mm]
\T_\pi(p\hRelease_{\precsim d}q)
  &= R\land\Next(\overline{Q}\WeakUntil D),\\[1mm]
\T_\pi(p\hRelease_{\succsim d}q)
  &= \Next(K\WeakUntil \Lambda).
\end{align*}
The ordinary timed operators have no translation clauses of their own.  Their
Coq definitions unfold into Boolean combinations containing the corresponding
hatted operator, after which the homomorphic Boolean translation and the
clauses above apply automatically.

All primitive formulas start logical evaluation after the current
position, which accounts for the $\Next$ in every clause.  The current timestamp
remains the clock reference, as required by the hatted semantics.

\subsection{Auxiliary Lemmas}
The proof uses elementary facts about LTL (the semantics of $\Always\Eventually$,
Eq.~\eqref{eq:sem-GF}, and the unfolding of Until and weak until) and clock
arithmetic: the value of a preserved clock accumulates the elapsed time
(Lemma~\ref{lem:clock-accumulation}), a reset clock measures the time elapsed
since its last reset, and the clock tests therefore bound the elapsed time
from above and from below (Eqs.~\eqref{eq:upper-clock}
and~\eqref{eq:lower-clock}), with the strict variants of
Section~\ref{subsec:strict-bounds}.  These lemmas and their proofs are given in
Section~\ref{sec:supp-lemmas}.

\subsection{Residual Obligations and the Canonical Extension}
\label{subsec:residuals}

Soundness can reason about an arbitrary clock-consistent extension.  Completeness,
however, must construct a reset trace that realizes the translation.  Following
the canonical extended-word construction of CASAAL~\cite{li_spin7}, the proof
does this using a deterministic canonical policy and residual predicates,
adapted here to formula-keyed shared clocks.

For upper-bounded Until, define
\begin{align*}
&\mathsf{URes}(i,v,d,p,q)\iff{}\\
&\quad\exists j\ge i.\;
  v+\elapsed_w(i,j)\le d
  \land w,j\satM q\\
&\qquad\land
  \forall k\in[i,j).\;w,k\satM p.
\end{align*}
The value $v$ represents time accumulated before the current position.  Thus
$\mathsf{URes}$ means that an older upper-bounded Until obligation is still able to
reach a valid $q$ endpoint.

For lower-bounded Release, define
\begin{align*}
&\mathsf{RRes}(i,v,d,p,q)\iff{}\\
&\quad\forall j\ge i.\;
  d\le v+\elapsed_w(i,j)
  \Longrightarrow{}\\
&\qquad\bigl(
  w,j\satM q
  \lor
  \exists k\in[i,j).\;w,k\satM p
\bigr).
\end{align*}
Here the residual represents a Release obligation that started before $i$ and has
already accumulated clock age $v$.

For the strict primitives, the proof uses the strict counterparts
\begin{align*}
&\mathsf{URes}^{<}(i,v,d,p,q)\iff{}\\
&\quad\exists j\ge i.\;
  v+\elapsed_w(i,j)< d
  \land w,j\satM q\\
&\qquad\land
  \forall k\in[i,j).\;w,k\satM p,\\[1mm]
&\mathsf{RRes}^{>}(i,v,d,p,q)\iff{}\\
&\quad\forall j\ge i.\;
  d< v+\elapsed_w(i,j)
  \Longrightarrow{}\\
&\qquad\bigl(
  w,j\satM q
  \lor
  \exists k\in[i,j).\;w,k\satM p
\bigr).
\end{align*}
They are the Coq predicates \texttt{ult\_\allowbreak{}residual} and
\texttt{rgt\_\allowbreak{}residual}; the one-step update of the latter is
\texttt{rgt\_\allowbreak{}residual\_\allowbreak{}shift}.

The reset decision is selected by matching on the \emph{formula-keyed clock}
itself, so the canonical policy has one case per hatted primitive; the strict
cases use the strict residual predicates $\mathsf{URes}^{<}$ and
$\mathsf{RRes}^{>}$ defined above.

\begin{definition}[Canonical reset policy and canonical extension]
\label{def:canonical-policy}
For a clock $x=F$, a position $i$, and a current clock value $v$, the reset
decision $\mathsf{reset}^\star(F,i,v)$ is defined as follows:
\begin{itemize}[leftmargin=*]
  \item $F=p\hUntil_{\le d}q$: the clock is \emph{preserved} (not reset) iff
        $v\le d$, $\mathsf{URes}(i,v,d,p,q)$, and $w,i\not\satM q$;
  \item $F=p\hUntil_{<d}q$: the clock is \emph{preserved} iff
        $v<d$, $\mathsf{URes}^{<}(i,v,d,p,q)$, and $w,i\not\satM q$;
  \item $F=p\hUntil_{\succsim d}q$ or $F=p\hRelease_{\precsim d}q$: the clock
        is reset iff $w,i\satM F$;
  \item $F=p\hRelease_{\ge d}q$: the clock is \emph{preserved} iff
        $v<d$, $\mathsf{RRes}(i,v,d,p,q)$, and $w,i\not\satM p$;
  \item $F=p\hRelease_{>d}q$: the clock is \emph{preserved} iff
        $v\le d$, $\mathsf{RRes}^{>}(i,v,d,p,q)$, and $w,i\not\satM p$.
\end{itemize}
The canonical extension $\rho^\star=\mathsf{canonical\_ext}(\mathit{root},w)$
has base word $w$ and, writing $u_i=\CVal_{\rho^\star}(i,x)$, clock values
\[
  u_0=0,\qquad
  u_{i+1}=
  \begin{cases}
    \Delta_i, & \text{if }\mathsf{reset}^\star(x,i,u_i),\\
    u_i+\Delta_i, & \text{otherwise},
  \end{cases}
\]
and reset decisions $\CReset_{\rho^\star}(i,x)=\mathsf{reset}^\star(x,i,u_i)$.
\end{definition}

This is the Coq definition \texttt{reset\_\allowbreak{}policy} with
\texttt{canonical\_\allowbreak{}val}.  The decisions depend on the future
of $w$ (through the residual predicates and the satisfaction of $F$),
so $\rho^\star$ is a non-constructive witness defined with classical logic; it
is used only to prove completeness.

The side conditions $w,i\not\satM q$ and $w,i\not\satM p$ are needed for the
canonical policy.  If the
older $\hUntil_{\precsim d}$ obligation is discharged by $q$ at the current position
$i$, nothing is required of the clock after $i$, and the clock is reset for the
newer activations; preserving it there could make a newer obligation miss its
bound.  Symmetrically, a $p$ at position $i$ discharges an older
$\hRelease_{\succsim d}$ obligation.  For $\hUntil_{\succsim d}$ and $\hRelease_{\precsim d}$,
restarting at every position where $F$ holds keeps the newest reference, and
the $\Gamma$ and weak-until continuations of the translation keep older
obligations represented.

The policy depends only on the clock $x=F$ and on $w$, not on the syntax-tree
path; this is why one decision simultaneously governs all syntax-tree
occurrences of the same timed formula.  Ordinary timed operators require no
additional policy cases because they unfold to these hatted primitives.
By construction,
\begin{equation}
  \mathsf{same\_base}(\rho^\star,w)
  \qquad\text{and}\qquad
  \mathsf{clock\_consistent}(\rho^\star).
\label{eq:canonical-properties}
\end{equation}

The residual predicates above are the mechanized form of the clock-reuse
principle stated informally in Section~\ref{subsec:clock-reuse}.  They allow the canonical policy to carry an older
reference forward without allocating a fresh clock for each activation.

\subsection{Soundness of the Primitive Timed Translations}
\label{subsec:timed-soundness}

For each primitive hatted timed constructor $F$, the local soundness lemma has
the form
\begin{equation}
  \rho,i\satL \T_\pi(F)
  \quad\Longrightarrow\quad
  \ewbase(\rho),i\satM F,
\label{eq:local-soundness}
\end{equation}
assuming clock consistency and the induction hypotheses for the immediate
subformulas.  No canonical reset-policy assumption appears in
\eqref{eq:local-soundness}.

\runinhead{Upper-bounded hatted Until.}
The translated LTL Until starts at $i+1$ and provides an endpoint $j>i$.
At position $i$ itself the translation does not constrain the clock, but the
recurrence~\eqref{eq:supp-clock-update} together with nonnegativity gives
$\CVal_\rho(i+1,x)\ge\Delta_i$ whether or not $x$ is reset at $i$.  Along the
intermediate positions $(i,j)$ the shared clock is unchanged, so
Lemma~\ref{lem:clock-accumulation} applied on $[i+1,j)$ yields
$\CVal_\rho(j,x)\ge\Delta_i+\elapsed_w(i+1,j)=\elapsed_w(i,j)$, and the
endpoint guard $x\le d$ gives
\[
  \elapsed_w(i,j)\le d.
\]
The induction hypotheses turn the translated prefix $\overline{P}$ into $p$ and the
endpoint $\overline{Q}$ into $q$, matching the semantics of
$p\hUntil_{\le d}q$.  For $\hUntil_{<d}$, the strict guard $x<d$ gives
$\elapsed_w(i,j)<d$ in the same way (\texttt{sound\_\allowbreak{}UhatLt}).

\runinhead{Lower-bounded hatted Until.}
The translation resets $x$ at $i$ and follows $\Gamma$.  In the finite branch,
a future state satisfies $x\ge d$ together with $\overline{P}\Until \overline{Q}$.  By
\eqref{eq:lower-clock}, the endpoint is at least $d$ time units after $i$, and
the nested Until supplies the required $p$-prefix and $q$-endpoint.  In the
second branch, $\Always \overline{P}\land\Always\Eventually \overline{Q}$ holds.  Then $p$ holds
forever and $q$ recurs infinitely often; time divergence guarantees a $q$
occurrence at physical distance at least $d$.  For $\hUntil_{>d}$, the strict
threshold $x>d$ gives a strict lower bound in the finite branch, and the
fairness branch uses the strict argument of Section~\ref{subsec:strict-bounds}
(\texttt{sound\_\allowbreak{}UhatGt}).

\runinhead{Upper-bounded hatted Release.}
The translation resets the shared clock and requires
\[
  \overline{Q}\WeakUntil\bigl((x>d)\lor(\overline{P}\land \overline{Q})\bigr).
\]
Before the weak-until endpoint, $q$ holds at every strictly future position.
If the endpoint is $\overline{P}\land \overline{Q}$, $p$ discharges the Release obligation; if it is
$x>d$, the upper-bounded window has ended.  For $\hRelease_{<d}$, the window
ends at $x\ge d$ (\texttt{sound\_\allowbreak{}RhatLt}).

\runinhead{Lower-bounded hatted Release.}
The translation requires
\[
  ((x<d)\land\Unch(x))
  \WeakUntil
  \bigl((\overline{P}\Release \overline{Q})\lor((x<d)\land \overline{P})\bigr)
\]
from $i+1$.  While the left side holds, the clock is preserved and remains
below $d$, so those positions are still before the lower threshold.  At the
endpoint either an ordinary untimed Release takes over or $p$ discharges the
obligation before the threshold is reached.  For $\hRelease_{>d}$, the
pre-threshold test is $x\le d$ (\texttt{sound\_\allowbreak{}RhatGt}).

There are no separate ordinary timed soundness lemmas.  After unfolding
Eqs.~\eqref{main-eq:derive-ule}--\eqref{main-eq:derive-rge}, their
correctness follows from the Boolean cases of the structural induction and the
eight hatted lemmas above.

\subsection{Completeness under the Canonical Reset Policy}
\label{subsec:timed-completeness}

The converse local result is proved for the canonical extension:
\[
  w,i\satM F
  \quad\Longrightarrow\quad
  \rho^\star,i\satL\T_\pi(F).
\]
The difficult point is that all occurrences of the same primitive timed formula
share one clock and each occurrence may itself be activated repeatedly.  Thus
several semantic obligations can compete for the same clock trace.  The
canonical reset policy and the residual predicates are designed to resolve
this interaction.

\runinhead{Upper-bounded hatted Until.}
Choose the first $q$-position among the MTL witnesses.  The proof maintains,
for every position $n$ with $i<n\le j$ up to that endpoint $j$, the invariant
\[
  \CVal_{\rho^\star}(n,x)+\elapsed_w(n,j)\le d.
\]
It need not hold at the activation position $i$ itself (the policy may reset
there because the old value is useless); it is established at $i+1$ in both
the reset and the preservation case.
If the policy preserves the clock, Definition~\ref{def:canonical-policy}
guarantees an older reference whose residual $\mathsf{URes}$ is still
satisfiable strictly after the current position; if it resets, the new value is
consistent with the current activation.  Hence every intermediate position satisfies the clock
guard required by the translated Until and the first $q$-position supplies its
endpoint.  This is the content of the mechanized lemma
\texttt{complete\_\allowbreak{}UhatLe}.  The strict case maintains the strict
invariant $\CVal_{\rho^\star}(n,x)+\elapsed_w(n,j)<d$ with the residual
$\mathsf{URes}^{<}$ (\texttt{complete\_\allowbreak{}UhatLt}).

\runinhead{Lower-bounded hatted Until.}
At every position where $F$ holds, the canonical policy resets the shared
clock (Definition~\ref{def:canonical-policy}).  The proof propagates an active obligation until a position satisfies
\[
  x\ge d
  \quad\text{and}\quad
  \overline{P}\Until \overline{Q},
\]
which yields the finite branch of $\Gamma$.  If no such position is reached,
the active-obligation invariant implies that $p$ remains continuously true,
and since every later position again carries a pending obligation that
requires a later $q$, the proposition $q$ recurs infinitely often.  This gives
\[
  \Always \overline{P}\land\Always\Eventually \overline{Q},
\]
the second branch of $\Gamma$.  This is formalized by
\texttt{complete\_\allowbreak{}UhatGe}, and by
\texttt{complete\_\allowbreak{}UhatGt} for the strict threshold $x>d$.

\runinhead{Upper-bounded hatted Release.}
The canonical policy resets the clock at every position where $F$ holds.  The
proof follows the currently active nontrivial Release obligation.  Until it
is discharged, the clock measures elapsed time from the relevant start and
$q$ continues to hold.  The weak until terminates either because $p\land q$
discharges the Release or because $x>d$, meaning that the upper-bounded window
has ended.  This is the role of \texttt{complete\_\allowbreak{}RhatLe}, and of
\texttt{complete\_\allowbreak{}RhatLt} for the strict window end $x\ge d$.

\runinhead{Lower-bounded hatted Release.}
Here the invariant is
\[
  \mathsf{RRes}
  \bigl(n,\CVal_{\rho^\star}(n,x),d,p,q\bigr).
\]
By Definition~\ref{def:canonical-policy}, the shared clock is preserved exactly
while $x<d$, the residual holds, and $p$ does not occur; in that case the
residual shifts
from $(i,v)$ to
\[
  (i+1,v+\Delta_i).
\]
The Coq lemma \texttt{rge\_\allowbreak{}residual\_\allowbreak{}shift} establishes this one-step update.
A converse lemma on the preservation policy recovers the residual information needed
at the next position, so the left side of the weak until remains valid until
either an ordinary untimed Release is established or $p$ occurs while the
clock is still below $d$.  This yields \texttt{complete\_\allowbreak{}RhatGe}.
The strict case uses $\mathsf{RRes}^{>}$, the pre-threshold test $x\le d$, and
the shift lemma \texttt{rgt\_\allowbreak{}residual\_\allowbreak{}shift}
(\texttt{complete\_\allowbreak{}RhatGt}).

As with soundness, ordinary timed operators need no additional completeness
lemmas: unfolding their definitions reduces them to Boolean structure and one
of the eight hatted primitive cases above.

\subsection{Structural Equivalence at Every Syntactic Occurrence}

The local timed lemmas are assembled with the homomorphic Boolean and untimed temporal
cases into the following stronger induction theorem.

\begin{theorem}[Canonical correctness at every occurrence]
\label{thm:canonical-all}
Let $\mathit{root}$ be well formed and suppose
\[
  \Node(\mathit{root},\pi)=f.
\]
Then for every position $i$,
\[
  w,i\satM f
  \quad\Longleftrightarrow\quad
  \rho^\star,i\satL\T_\pi(f),
\]
where $\rho^\star=\mathsf{canonical\_ext}(\mathit{root},w)$.
\end{theorem}

\begin{proof}[Proof idea]
The proof is by structural induction on the formula $f$.  Atomic propositions,
Boolean connectives, Next, untimed Until, and untimed Release follow directly
from their semantics and the induction hypotheses.  Since the ordinary timed
operators are Coq definitions, they have already unfolded into these structural
cases.  The only primitive timed branches are therefore the eight hatted
constructors.  In each such branch, the forward direction invokes the
corresponding completeness lemma from
Section~\ref{subsec:timed-completeness}, while the reverse direction invokes
the corresponding soundness lemma from
Section~\ref{subsec:timed-soundness}.  The $\Node$ lemmas locate the two
children, and well-formedness of the located formula provides the bound and
recursive hypotheses.
\end{proof}

This is stronger than a root-only statement and is important in the mechanized
proof.  The induction retains the path in order to locate the current subformula
and its children, but the path is not the clock key.  If two paths locate the
same primitive timed formula, both branches refer to the same formula-keyed
clock.

\subsection{Proof of the Correctness Theorem}
\label{subsec:proof-main-theorem}
We now prove Theorem~\ref{main-thm:encoding-correctness}.

\begin{proof}
We prove the two implications separately.

\smallskip
\noindent\emph{($\Rightarrow$) MTL to clocked LTL.}
Assume $w,0\satM f$.  Choose the canonical extension
\[
  \rho=\mathsf{canonical\_ext}(f,w).
\]
By \eqref{eq:canonical-properties}, $\rho$ has the same base timed word as $w$ and is
clock-consistent.  Applying Theorem~\ref{thm:canonical-all} at the root path $[]$ and
position $0$ gives
\[
  \rho,0\satL\T(f).
\]
Therefore the required extension exists.

\smallskip
\noindent\emph{($\Leftarrow$) Clocked LTL to MTL.}
Assume there exists an arbitrary $\rho$ such that
\[
\begin{gathered}
  \mathsf{same\_base}(\rho,w),\qquad
  \mathsf{clock\_consistent}(\rho),\\
  \rho,0\satL\T(f).
\end{gathered}
\]
This extension need not be canonical, so completeness cannot be used.  Instead, a
structural induction on $f$ uses the policy-independent soundness lemmas of
Section~\ref{subsec:timed-soundness} and obtains
\[
  \ewbase(\rho),0\satM f.
\]
Finally, $\mathsf{same\_base}(\rho,w)$ rewrites the underlying timed word
$\ewbase(\rho)$ to $w$, yielding $w,0\satM f$.
\end{proof}

\runinhead{Scope of the mechanized result.}
Theorem~\ref{main-thm:encoding-correctness} covers the MTL formula, its
clocked-LTL encoding, and clock-consistent extensions. Corollary
\ref{main-cor:mtl-tba-correct} extends the result, through relaxation and reset
completion, to the timed automaton built from any backend automaton satisfying
the stated propositional-language contract. Initial clock values remain
existential in this theorem. Agreement with conventional zero initialization,
the parser and external backend interface, and output printing are outside
the result proved here.
